% open-logic.tex
% compiles to produce complete text using open-logic-dev style

\documentclass[../include/open-logic-part]{subfiles}

\begin{document}

\clearpage

\begin{editorial}
ही फाइल ओपन लॉजिक प्रकल्पातील सर्व सामग्री लोड करते. अशा
संपादकीय नोंदी दिसत असल्यास, ही सामग्री कशी सादर केली जाणार
याचा विचार न करता फाइल संकलित केली गेली आहे, हे त्यावरून
समजते. तुम्हाला हे वाचता येत असेल, तर या PDF वरून शिकवणे
किंवा अभ्यास करणे बहुधा \emph{योग्य नाही}.

अध्यापनासाठी किंवा स्व-अध्ययनासाठी अधिक योग्य मजकूर तयार
करता यावा यासाठी ओपन लॉजिक प्रकल्पात अनेक सोयी आहेत.
उदाहरणार्थ, नेहमीच्या मांडणीत सर्व तार्किक संयोजक आदिम
मानले जातात आणि पुराव्यांमध्ये प्रत्येक संयोजकाचे सर्व प्रकार
तपासले जातात. पण त्यांपैकी काही प्रकार सरावासाठी सोडणे अधिक
चांगले असते. ओपन लॉजिक प्रकल्पाचे कामही अजून सुरू आहे.
सहयोग आणि सुधारणा वाढाव्यात म्हणून मसुदा स्वरूपातील किंवा
सरावप्रश्न नसलेली सामग्रीसुद्धा येथे समाविष्ट आहे. अभ्यासक्रमासाठी
तयार केलेल्या PDF मध्ये असे विभाग वगळले जातात.

अध्यापन आणि अभ्यासासाठी अधिक योग्य PDF पाहण्यासाठी
\href{http://builds.openlogicproject.org/}{OLP संकेतस्थळावरील नमुना अभ्यासक्रम}
पहा. स्वतःचा मजकूर तयार करण्यासाठी
\href{https://github.com/OpenLogicProject/OpenLogic/tree/master/courses/sample}{नमुना चालक फाइल}
वापरून सुरुवात करता येईल; अधिक सजवलेल्या आणि प्रगत उदाहरणांसाठी
या प्रकल्पातून तयार झालेल्या पाठ्यपुस्तकांचे स्रोत पाहता येतील.
\end{editorial}

\olimport[sets-functions-relations]{sets-functions-relations-complete}

\olimport[propositional-logic]{propositional-logic}

\olimport[first-order-logic]{first-order-logic}

\olimport[model-theory]{model-theory}

\olimport[computability]{computability}

\olimport[turing-machines]{turing-machines}

\olimport[incompleteness]{incompleteness}

\olimport[second-order-logic]{second-order-logic}

\olimport[lambda-calculus]{lambda-calculus}

\olimport[many-valued-logic]{many-valued-logic}

\olimport[normal-modal-logic]{normal-modal-logic}

\olimport[applied-modal-logic]{applied-modal-logic}

\olimport[intuitionistic-logic]{intuitionistic-logic}

\olimport[counterfactuals]{counterfactuals}

\olimport[set-theory]{set-theory}

\olimport[methods]{methods}

\olimport[history]{history}

\olimport[reference]{reference}

\end{document}
